% source: https://github.com/pfradin/luadraw/discussions/312#discussioncomment-18080177 (pfradin, block 1)
\documentclass{standalone}
\usepackage[3d]{luadraw}
\usepackage[svgnames]{xcolor}
\usepackage{fourier-otf}
\begin{document}
\begin{luadraw}{name=3dobjects}
local ld = luadraw
local pt3d, cpx = ld.pt3d, ld.cpx
local M, Mc, Ms = pt3d.M, pt3d.Mc, pt3d.Ms
local a = 5 --angle in degrees
local Origin,vecI,vecJ,vecK = pt3d.Origin, pt3d.vecI, pt3d.vecJ, pt3d.vecK
local g = ld.graph3d:new{window = {-1.5,1.5,-0.5,2.5}, size={12,12}, viewdir = {-a,90}}
ld.Hiddenlinestyle = "dashed"; ld.Hiddenlines = true
g:Dsphere(vecK,1,{color="violet!50", mode=ld.mBorder})  --no equator
local N = vecJ
local A, B = table.unpack( g:Sphere_tangency(vecK,1,{vecK, N}) )
g:Dpath3d({A,vecK,B,1,-1,N,"ca"})
g:Dpath3d({A,vecK,B,1,1,N,"ca"},"dashed")
local N = Ms(1,-a*ld.deg,a*ld.deg)
local A, B = table.unpack( g:Sphere_tangency(vecK,1,{vecK, N}) )
g:Dpath3d({A,vecK,B,1,-1,N,"ca"})
g:Dpath3d({A,vecK,B,1,1,N,"ca"},"dashed")
g:Dline3d({vecK,vecK}, 'red')
g:Dline3d({vecK,vecJ}, 'red')
g:Show()
\end{luadraw}
\end{document}
